Линейность матожидания: $M(x_1 + x_2 + \dots + x_n) = M(x_1) + M(x_2) + \dots + M(x_n)$
\subsection{Sequence sums}
Сумма геометрической прогрессии: $S_n = \frac{b_1 (1 - q^n)}{1 - q}$ \\
Сумма бесконечно убывающей геометрической прогрессии: $S = \frac{b_1}{1 - q}$ \\
Сумма арифметической прогрессии: $S_n = \frac{a_1 + a_n}{2} \cdot n = \frac{2a_1 + d(n - 1)}{2} \cdot n$ \\
$1^2 + 2^2 + 3^2 + \dots + n^2 = \frac{n(n + 1)(2n + 1)}{6}$ \\
$1^3 + 2^3 + 3^3 + \dots + n^3 = \frac{n^2 (n + 1)^2}{4}$ \\
\subsection{Calculus}
Площадь криволинейного сектора:\\
$\mu(F) = \frac{1}{2} \int_{\alpha}^{\beta} r^2(\varphi) d\varphi$ \\\\
Интегралы:
\begin{enumerate}
	\item $ \int \sqrt{a^2 - x^2} dx = \frac{x}{2} \sqrt{a^2 - x^2} + \frac{a^2}{2} arcsin \frac{x}{a} + C $ \\
	\item $ \int \sqrt{x^2 \pm a^2} dx = \frac{x}{2} \sqrt {x^2 \pm a^2} \pm \frac{a^2}{2} ln |x + \sqrt{x^2 \pm a^2}| + C $ \\
	\item $ \int \frac{dx}{a^2 + x^2} = \frac{1}{a} arctg \frac{x}{a} + C (a \ne 0) $ \\
	\item $ \int \frac{dx}{a^2 - x^2} = \frac{1}{2a} ln |\frac{a + x}{a - x}| + C (a \ne 0) $ \\
	\item $ \int \frac{x dx}{a^2 \pm x^2} = \pm \frac{1}{2} ln |a^2 \pm x^2| + C $ \\
	\item $ \int \frac{dx}{\sqrt{a^2 - x^2}} = arcsin \frac{x}{a} + C (a \ne 0) $ \\
	\item $ \int \frac{dx}{\sqrt{x^2 \pm a^2}} = ln |x + \sqrt{x^2 \pm a^2}| + C $ \\
	\item $ \int \frac{x dx}{\sqrt{a^2 \pm x^2}} = \pm \sqrt{a^2 \pm x^2} + C $ \\
\end{enumerate}
\subsection{Combinatorics}
\textbf{Размещения}: 
\begin{itemize}
	\item Без повторений: $A_n^k = \frac{n!}{(n - k)!} $
	\item С повторениями: $\overline{A}_n^k = n^k $
\end{itemize}
\textbf{Сочетания}:
\begin{itemize}
	\item Без повторений: $C_n^k = \frac{n!}{k! (n - k)!}$
	\item С повторениями: $\overline{C}_n^k = \frac{(n + k - 1)!}{k! (n - 1)!} $
\end{itemize}
\textbf{Перестановки}:
\begin{itemize}
	\item Без повторений: $P_n = n!$
	\item С повторениями: $P_n(n_1, n_2, \dots, n_k) = \frac{n!}{n_1! n_2! \dots n_k!} $
\end{itemize}
\textbf{Биномиальные коэффициенты}
\begin{enumerate}
	\item $C_n^k = C_{n - 1}^k + C_{n - 1}^{k - 1} $;
	\item $C_n^k = C_{n}^{n - k} $;
	\item $\binom{n}{k}=\frac nk \binom{n-1}{k-1}$;
	\item $\binom{n}{k}$ можно получить из $\binom{n}{0}$ за $O(k)$: $\binom{n}{k}=\frac{n-k+1}k \binom{n}{k-1}$;
	\item $A_n^k=\frac{n!}{(n-k)!}=C_n^k\:k!$ --- количество размещений из $n$ по $k$;
	\item $\sum\limits_{k=0}^{n}k C_n^k=n2^{n-1}$, $\sum\limits_{k=0}^{n}k^2 C_n^k=(n+n^2)2^{n-2}$;
	\item $\sum\limits_{k=0}^{n}{(C_n^{k})^2}=C_{2n}^n$;
	\item $\sum\limits_{k=0}^{\lfloor\frac{n}{2}\rfloor}C_{n-k}^k=F(n+1)$;
	\item $\sum\limits_{k\le n} \binom{r+k}{k}=\binom{r+n+1}{n}, n\in \mathbb{Z}$;
	\item $\sum\limits_{0\le k\le n} \binom{k}{m} = \binom{n+1}{m+1}, m, n\in \mathbb{Z}, m, n\ge 0$;
	\item $\sum\limits_k \binom{r}{k}\binom{s}{n-k}=\binom{r+s}{n}, n\in \mathbb{Z}$.
\end{enumerate}
\textbf{Числа Каталана}:
1, 1, 2, 5, 14, 42, 132, 429, 1430, 4862, 16796...

$C_0 = 1, C_n = \sum_{i=0}^{n-1}C_i C_{n-1-i}$, для $n \ge 1$.

$C_n = \frac{1}{n+1}\binom{2n}{n} = \binom{2n}{n} - \binom{2n}{n - 1}$.

Применения:
\begin{itemize}
	\item Правильные скобочные последовательности из n пар скобок;
	\item Количество бинарных деревьев с n листьями;
	\item Количество неизоморфных деревьев из n вершин;
	\item Триангуляции n"=угольника;
	\item Таблицы Юнга размера 2xn (в каждой строке числа возрастают).
\end{itemize}
\textbf{Stars and bars}:

Количество способов представить число n в виде суммы k неотрицательных слагаемых: $C_{n + k - 1}^{k - 1} $

Количество способов представить число n в виде суммы k положительных слагаемых: $C_{n - 1}^{k - 1} $

Для слагаемых больше 1: $C_{n - k - 1}^{k - 1}$. 

Сумма произведений всех подмножеств множества a: $(\prod_{i=1}^{n} (a_i + 1)) - 1$ (аналогично формуле Эйлера для подсчёта количества делителей числа)

\textbf{Числа Стирлинга:}


\textbf{Числа I рода}.

Без знака $c(n,k)$: количество перестановок из $n$ элементов с $k$ циклами.

Со знаком $s(n,k)$: $(x)_n = \sum_{k=0}^n s(n,k)x^k$, $(x)_n = x(x-1)(x-2)\dots(x-n+1)$.

$c(n, k) = |s(n,k)| = (-1)^{n-k}s(n,k)$.

$s(0,0)=c(0,0)=1$.

$s(n,0)=c(n,0)=0$, $s(0,k)=c(0,k)=0$.

$s(n,k) = s(n-1,k-1) - (n-1) \cdot s(n-1,k)$.

$c(n,k) = c(n-1,k-1) + (n-1) \cdot c(n-1,k)$.

\begin{center}
    \begin{tabular}{|l|l|l|l|l|l|l|l|} \hline
        n | k      & \textbf{0} & \textbf{1} & \textbf{2} & \textbf{3} & \textbf{4} & \textbf{5} & \textbf{6} \\ \hline
        \textbf{0} & 1          &            &            &            &            &            &            \\ \hline
        \textbf{1} & 0          & 1          &            &            &            &            &            \\ \hline
        \textbf{2} & 0          & -1         & 1          &            &            &            &            \\ \hline
        \textbf{3} & 0          & 2          & -3         & 1          &            &            &            \\ \hline
        \textbf{4} & 0          & -6         & 11         & -6         & 1          &            &            \\ \hline
        \textbf{5} & 0          & 24         & -50        & 35         & -10        & 1          &            \\ \hline
        \textbf{6} & 0          & -120       & 274        & -225       & 85         & -15        & 1          \\ \hline
    \end{tabular}
\end{center}

\textbf{Числа II рода}.

$S(n,k)$ --- количество неупорядоченных разбиений $n$-элементного множества на $k$ непустых подмножеств.

$S(n,k) = S(n - 1, k - 1) + k \cdot S(n-1,k)$ для $0 < k \le n$.

$S(n,k) = \sum_{j=0}^{n-1} \binom{n-1}{j}S(j, k-1)$.

$S(0,0) = 1, S(n,0)=0, S(j,k)=0$ при $k>j$.

$S(n, k) = \frac{1}{k!} \sum_{j=0}^k (-1)^{k+j} \binom{k}{j} j ^ n$.

\begin{center}
    \begin{tabular}{|l|l|l|l|l|l|l|l|l|} \hline
        n | k      & \textbf{0} & \textbf{1} & \textbf{2} & \textbf{3} & \textbf{4} & \textbf{5} & \textbf{6} & \textbf{7} \\ \hline
        \textbf{0} & 1          &            &            &            &            &            &            &            \\ \hline
        \textbf{1} & 0          & 1          &            &            &            &            &            &            \\ \hline
        \textbf{2} & 0          & 1          & 1          &            &            &            &            &            \\ \hline
        \textbf{3} & 0          & 1          & 3          & 1          &            &            &            &            \\ \hline
        \textbf{4} & 0          & 1          & 7          & 6          & 1          &            &            &            \\ \hline
        \textbf{5} & 0          & 1          & 15         & 25         & 10         & 1          &            &            \\ \hline
        \textbf{6} & 0          & 1          & 31         & 90         & 65         & 15         & 1          &            \\ \hline
        \textbf{7} & 0          & 1          & 63         & 301        & 350        & 140        & 21         & 1          \\ \hline
    \end{tabular}
\end{center}

1-го рода: $Z_n^k =(n-1)Z_{n - 1}^k + Z_{n-1}^{k-1}$ --- число способов разбиения множества из $n$ элементов на $k$ циклов (циклы нельзя переворачивать, только прокручивать).

2-го рода: $S_n^k = kS_{n - 1}^k + S_{n-1}^{k-1} = \frac1{k!} \sum\limits_{j=0}^k (-1)^{k+j}\binom{k}{j}j^n$ --- число способов разбиения множества из $n$ элементов на $k$ непустых подмножеств.

\textbf{Количество неубывающих последовательностей из $n$ элементов от $0$ до $a$ равно $\binom{n + a}{n}$.}

\textbf{Перестановки без неподвижных точек:} $cnt = n! \sum\limits_{i = 0}^n \frac{(-1)^n}{i!} \approx \frac{n!}e$



\textbf{Числа Белла}

Число всех неупорядоченных разбиений $n$-элементного множества.

1, 1, 2, 5, 15, 52, 203, 877, 4'140, 21'147, 115'975, 678'570, 4'213'597, \dots

Способы распределить $n$ пронумерованных шаров по $n$ идентичным коробкам. Способы разложить на множители составное число из $n$ простых множителей.

$B_n = \sum_{m=0}^{n}S(n,m)$ (сумма чисел Стирлинга II рода).

$B_{n+1} = \sum_{k=0}^n \binom{n}{k}B_k$

Для простого $p$:

$B_{n+p} \equiv B_n + B_{n+1} (mod\ p)$
